\begingroup
\renewcommand{\qD}{q_D}
\section{Foundation F129: the exact disk-source envelope}
\label{module:F129}
\noindent\textit{Audited source: \nolinkurl{convex_largeN_polya_szego_stage129_EF_sharp_strip_gluing_closure.tex}.}
\subsection{Exact source and the three disk rows}

Let
\begin{equation}\label{F129:eq:fan}
 \alpha_i>0,\qquad \sum_i\alpha_i=2\pi,\qquad
 a=\frac{2\pi}{N},\qquad
 \Jfour(\alpha)=\sum_i\alpha_i^4-Na^4.
\end{equation}
Choose consecutive side normals \(\theta_i\), put
\(I_i=[\theta_i,\theta_{i+1}]\), and use \(x=\theta-\theta_i\) on
\(I_i\).  After removal of a constant, the unnormalized exact secant source
is
\begin{equation}\label{F129:eq:B}
 B_\alpha(\theta_i+x)
 =F_{\alpha_i}(x)
 :=\sec\bigl(\min\{x,\alpha_i-x\}\bigr)-1,
 \qquad 0\le x\le\alpha_i.
\end{equation}
Both \(F_h\) and \(F_h'\) vanish at \(0,h\), so the pieces form one
continuous periodic source.  The exact area-normalized source is
\begin{equation}\label{F129:eq:b}
 b_\alpha=h_\alpha\PiNG B_\alpha,\qquad
 h_\alpha^2=\frac{\pi}{A_\alpha},\qquad
 A_\alpha=\sum_i\tan\frac{\alpha_i}{2}.
\end{equation}

For \(n\ge2\), set
\begin{equation}\label{F129:eq:symbols}
 k_n=j_{0,1}\frac{J_{n+1}(j_{0,1})}{J_n(j_{0,1})},
 \qquad m_n=n+1-k_n.
\end{equation}
On nongeometric modes,
\begin{equation}\label{F129:eq:split}
 \qD(f)=\qdot(f)+\qzero(f)-\kform(f),
\end{equation}
where
\begin{align}
 \qdot(f)&=\sum_{|n|\ge2}|n|\,|\widehat f_n|^2,\label{F129:eq:qdot}\\
 \qzero(f)&=\sum_{|n|\ge2}|\widehat f_n|^2,\label{F129:eq:qzero}\\
 \kform(f)&=\sum_{|n|\ge2}k_{|n|}|\widehat f_n|^2.\label{F129:eq:kform}
\end{align}
Stages 51 and 65 give, for \(0\le r\le3\),
\begin{equation}\label{F129:eq:k-symbol}
 0<k_n\le\frac Cn,\qquad
 |\Delta^rk_n|\le C_r(1+n)^{-1-r}.
\end{equation}

\subsection{The sharp broken-strip gluing lemma}

For \(0<h\le h_0<\pi\), let \(V_h\) be the Neumann harmonic extension of
\(F_h\) to the half-strip
\(\Sigma_h=(0,h)\times(0,\infty)\):
\begin{equation}\label{F129:eq:strip-problem}
 \Delta V_h=0,\qquad V_h(\cdot,0)=F_h,\qquad
 \partial_xV_h(0,\cdot)=\partial_xV_h(h,\cdot)=0.
\end{equation}
Its limit at infinity is the cell mean
\begin{equation}\label{F129:eq:mean}
 \mathfrak m(h)=\frac1h\int_0^hF_h(x)\,dx,
\end{equation}
and put
\begin{equation}\label{F129:eq:N}
 \mathcal N(h)=\int_{\Sigma_h}|\nabla V_h|^2\,dx\,dy.
\end{equation}

\begin{lemma}[Scaled strip data]\label{F129:lem:scaled}
For \(0<h\le h_0\),
\begin{equation}\label{F129:eq:scaled-bounds}
 |\mathfrak m(h)|\le Ch^2,\qquad
 0\le\mathcal N(h)\le Ch^4,\qquad
 |\mathcal N''(h)|\le Ch^2.
\end{equation}
If \(w_h(y)=V_h(0,y)=V_h(h,y)\), then
\begin{equation}\label{F129:eq:wall-scale}
 w_h(y)=h^2W_h(y/h),
\end{equation}
where \(h\mapsto W_h\) is bounded with two derivatives in
\(H^{1/2}(\mathbb R_+)\cap L^2(\mathbb R_+)\) after its constant limit is
removed.
\end{lemma}

\begin{proof}
Write
\begin{equation}\label{F129:eq:profile}
 F_h(hs)=h^2\Phi_h(s),\qquad
 \Phi_h(s)=\frac{\sec(h\min\{s,1-s\})-1}{h^2}.
\end{equation}
The family \(\Phi_h\), including its analytic value at \(h=0\), is bounded
with two \(h\)-derivatives and four one-sided \(s\)-derivatives.  Its value
and first derivative vanish at \(s=0,1\).  Expand it in the Neumann cosine
basis:
\begin{equation}
 \Phi_h(s)=\overline\Phi_h+\sum_{\ell\ge1}c_\ell(h)\cos(\ell\pi s).
\end{equation}
The midpoint derivative jump gives
\(|\partial_h^r c_\ell(h)|\le C_r(1+\ell)^{-2}\), \(0\le r\le2\).
Consequently
\begin{align}
 V_h(hs,hy)
 &=h^2\left\{\overline\Phi_h+
   \sum_{\ell\ge1}c_\ell(h)
   \cos(\ell\pi s)e^{-\ell\pi y}\right\},\label{F129:eq:scaled-V}\\
 \mathcal N(h)
 &=\frac{\pi h^4}{2}\sum_{\ell\ge1}\ell\,|c_\ell(h)|^2.\label{F129:eq:N-series}
\end{align}
Equations \eqref{F129:eq:scaled-bounds} and \eqref{F129:eq:wall-scale} now follow by
termwise differentiation; the series are dominated by
\(\sum\ell^{-3}\).
\end{proof}

For \(H>0\), use the wall-trace norm
\begin{equation}\label{F129:eq:XH}
 \|g\|_{X_H}^2
 :=\|g\|_{\dot H^{1/2}(\mathbb R_+)}^2
   +H^{-1}\|g\|_{L^2(\mathbb R_+)}^2.
\end{equation}
Functions below vanish at \(y=0\) and at infinity, so the homogeneous
half-line norm has no constant ambiguity.

\begin{lemma}[Wall equalization and lifting]\label{F129:lem:wall}
Fix \(a\le h_0\).  There are modified strip fields
\(\widetilde V_h\) with the same bottom value \(F_h\), common limit
\(\mathfrak m(a)\) at infinity, and
\begin{equation}\label{F129:eq:modified-energy}
 \int_{\Sigma_h}|\nabla\widetilde V_h|^2
 \le\mathcal N(h)+C\{\mathfrak m(h)-\mathfrak m(a)\}^2.
\end{equation}
For two adjacent widths \(h,k\), \(H=\max\{h,k\}\), their wall jump
\(j_{h,k}\) satisfies
\begin{equation}\label{F129:eq:jump}
 \|j_{h,k}\|_{X_H}^2
 \le C\{(h^2-a^2)^2+(k^2-a^2)^2\}.
\end{equation}
Moreover the jump can be lifted into the wider of the two adjacent strips,
with zero bottom value, zero limit at infinity, and Dirichlet energy at
most \(C\|j_{h,k}\|_{X_H}^2\).
\end{lemma}

\begin{proof}
Choose a fixed smooth \(\chi\) on \([0,\infty)\) with
\(\chi(0)=0\) and \(\chi=1\) on \([1,\infty)\), and set
\begin{equation}\label{F129:eq:tail-correction}
 C_h(x,y)=-\{\mathfrak m(h)-\mathfrak m(a)\}\chi(y/h),
 \qquad \widetilde V_h=V_h+C_h.
\end{equation}
The correction is independent of \(x\).  Since the horizontal mean of
\(\partial_yV_h\) is zero, its Dirichlet cross term with \(C_h\) vanishes.
Direct scaling gives
\begin{equation}
 \int_{\Sigma_h}|\nabla C_h|^2
 =|\mathfrak m(h)-\mathfrak m(a)|^2
   \int_0^\infty|\chi'(s)|^2\,ds,
\end{equation}
which proves \eqref{F129:eq:modified-energy}.

The modified wall trace is zero at \(y=0\) and tends to the common value
\(\mathfrak m(a)\).  Subtract that common value.  Equations
\eqref{F129:eq:wall-scale} and \eqref{F129:eq:tail-correction}, followed by scaling
in \eqref{F129:eq:XH}, prove \eqref{F129:eq:jump}.  For completeness, first suppose
\(h,k\in[a/2,2a]\).  The fundamental theorem of calculus in the scale
parameter gives
\[
 \|j_{h,k}\|_{X_H}^2\le C(h-k)^2(h+k)^2.
\]
If at least one width lies outside \([a/2,2a]\), the triangle inequality
and scaling of the two tail-equalized traces give
\[
 \|j_{h,k}\|_{X_H}^2
 \le C(h^4+k^4+a^4)
 \le C\{(h^2-a^2)^2+(k^2-a^2)^2\}.
\]
When one width remains in \([a/2,2a]\), insert the nearer endpoint
\(a/2\) or \(2a\) as an intermediate scale, use the first estimate on the
short piece, and the coarse estimate on the other.  This proves
\eqref{F129:eq:jump} in every ratio.  The \(a^4\) term is essential when both
\(h\) and \(k\) are much smaller than \(a\): it pays for changing their
common strip tail to the regular-cell mean.

Finally, the Poisson extension in a quadrant is a right inverse of the
vertical trace with energy
\(\|j\|_{\dot H^{1/2}(\mathbb R_+)}^2\).  Multiply it by a horizontal cutoff
supported before distance \(H/3\).  The cutoff row costs
\(CH^{-1}\|j\|_2^2\), proving the \(X_H\) estimate.  Put the lifting in the
wider adjacent strip.  The two possible liftings assigned to one strip
occupy its left and right thirds, so their overlap is uniformly bounded.
\end{proof}

\begin{proposition}[Sharp no-ratio strip gluing]\label{F129:prop:gluing}
For every sufficiently fine positive fan,
\begin{equation}\label{F129:eq:gluing}
 \sum_{n\ne0}|n|\,|\widehat B_{\alpha,n}|^2
 \le\frac1{2\pi}\sum_i\mathcal N(\alpha_i)
   +C\sum_i(\alpha_i^2-a^2)^2.
\end{equation}
At the regular fan equality holds without the last term:
\begin{equation}\label{F129:eq:uniform-strip}
 \sum_{n\ne0}|n|\,|\widehat B_{a\mathbf1,n}|^2
 =\frac N{2\pi}\mathcal N(a).
\end{equation}
\end{proposition}

\begin{proof}
Place \(\widetilde V_{\alpha_i}\) over \(I_i\).  Lemma~\ref{F129:lem:wall}
gives wall jumps which vanish at the bottom and at infinity.  Lift all
jumps as in that lemma, assigning each lifting to the wider adjacent strip.
The resulting field \(U\) belongs to
\(H^1(\Sone\times\mathbb R_+)\), has bottom trace \(B_\alpha\), and tends
to the single constant \(\mathfrak m(a)\).

The original \(V_{\alpha_i}\) is harmonic and has zero normal derivative
on its vertical walls.  Hence it is Dirichlet-orthogonal to every jump
lifting, which vanishes on the bottom and at infinity.  It is also
orthogonal to the constant-in-\(x\) tail correction.  The only remaining
cross term is between tail corrections and jump liftings; Cauchy--Schwarz
absorbs it into their two energies.  Lemma~\ref{F129:lem:wall} therefore gives
\begin{equation}
 \int|\nabla U|^2
 \le\sum_i\mathcal N(\alpha_i)
 +C\sum_i(\alpha_i^2-a^2)^2.
\end{equation}
The harmonic-extension characterization of the periodic
\(\dot H^{1/2}\) norm proves \eqref{F129:eq:gluing}.

For \(\alpha_i=a\), all strip fields and their wall traces agree.  No tail
correction or lifting is present, and reflection of \(V_a\) across every
vertical wall is exactly the periodic harmonic extension.  This proves
\eqref{F129:eq:uniform-strip}.
\end{proof}

\subsection{The order-one and identity source envelopes}

Recall the exact identity
\begin{equation}\label{F129:eq:Jidentity}
 \Jfour(\alpha)
 =\sum_i(\alpha_i-a)^2
   (\alpha_i^2+2a\alpha_i+3a^2).
\end{equation}
In particular,
\begin{equation}\label{F129:eq:square-charge}
 \sum_i(\alpha_i^2-a^2)^2
 \le\Jfour(\alpha).
\end{equation}

\begin{theorem}[Sharp order-one source envelope]\label{F129:thm:qdot}
For every sufficiently fine positive fan,
\begin{equation}\label{F129:eq:qdot-envelope}
 \qdot(\PiNG B_\alpha)-\qdot(\PiNG B_{a\mathbf1})
 \le C\Jfour(\alpha).
\end{equation}
\end{theorem}

\begin{proof}
Removing the modes \(\pm1\) only decreases the left energy at \(\alpha\);
they vanish at the regular fan by cyclic symmetry.  Apply
Proposition~\ref{F129:prop:gluing} and subtract \eqref{F129:eq:uniform-strip}.
By Lemma~\ref{F129:lem:scaled} and Taylor's formula,
\begin{align}
 \sum_i\{\mathcal N(\alpha_i)-\mathcal N(a)\}
 &=\sum_i\{\mathcal N(\alpha_i)-\mathcal N(a)
           -\mathcal N'(a)(\alpha_i-a)\}\notag\\
 &\le C\sum_i(\alpha_i-a)^2(\alpha_i+a)^2
 \le C\Jfour(\alpha).\label{F129:eq:N-Bregman}
\end{align}
Use \eqref{F129:eq:square-charge} for the gluing row.
\end{proof}

\begin{lemma}[Identity source envelope]\label{F129:lem:qzero}
For every sufficiently fine positive fan,
\begin{equation}\label{F129:eq:qzero-envelope}
 \qzero(\PiNG B_\alpha)-\qzero(\PiNG B_{a\mathbf1})
 \le C\Jfour(\alpha).
\end{equation}
\end{lemma}

\begin{proof}
Put
\begin{equation}
 R(h)=\int_0^hF_h(x)^2\,dx,\qquad
 M(h)=\int_0^hF_h(x)\,dx
     =2\int_0^{h/2}(\sec x-1)\,dx.
\end{equation}
Scaling gives \(|R''(h)|\le Ch^3\).  The linear terms cancel under
\(\sum_i(\alpha_i-a)=0\), so
\begin{equation}
 \sum_i\{R(\alpha_i)-R(a)\}
 \le C h_0\sum_i(\alpha_i-a)^2(\alpha_i+a)^2
 \le C\Jfour(\alpha).
\end{equation}
Moreover
\begin{equation}
 M''(h)=\frac12\sec\frac h2\tan\frac h2>0.
\end{equation}
Thus Jensen's inequality gives
\(\sum_iM(\alpha_i)\ge NM(a)>0\); subtracting the squared global mean can
only improve the relative upper bound.  Finally, the regular source has no
modes \(\pm1\), while deleting those modes at \(\alpha\) again decreases
the energy.  Parseval proves \eqref{F129:eq:qzero-envelope}.
\end{proof}

\subsection{The smoother Bessel row}

\begin{lemma}[Regular Bessel baseline]\label{F129:lem:kbaseline}
For every sufficiently small \(a\),
\begin{equation}\label{F129:eq:kbaseline}
 0\le\kform(\PiNG B_{a\mathbf1})\le Ca^5.
\end{equation}
\end{lemma}

\begin{proof}
Cyclic symmetry restricts the Fourier support to \(N\mathbb Z\).
The scaled cosine expansion in Lemma~\ref{F129:lem:scaled} gives
\(|\widehat B_{\ell N}|\le Ca^2(1+|\ell|)^{-2}\).  Since
\(k_{\ell N}\le C/(|\ell|N)\le Ca/|\ell|\), summation proves
\eqref{F129:eq:kbaseline}.
\end{proof}

\begin{lemma}[Quasiuniform Bessel material Hessian]\label{F129:lem:kHessian}
Suppose
\begin{equation}
 \frac a2\le\gamma_i\le\frac{3a}{2},\qquad
 \sum_i\gamma_i=2\pi,\qquad \sum_i d_i=0.
\end{equation}
Then
\begin{equation}\label{F129:eq:kHessian}
 \left|D^2\{\kform(\PiNG B_\gamma)\}[d,d]\right|
 \le Ca\sum_i\gamma_i^2d_i^2.
\end{equation}
\end{lemma}

\begin{proof}
Pull every normal interval affinely to the regular material interval.
The profile \eqref{F129:eq:profile} and its first two material derivatives are
uniformly bounded after the factors \(\gamma_i^2\), \(\gamma_id_i\), and
\(d_i^2\) are removed.  The multiplier bounds \eqref{F129:eq:k-symbol} allow
three cyclic Abel summations in the pulled-back kernel.  The first removes
the arbitrary root velocity, because consecutive midpoint velocities
differ by \((d_i+d_{i+1})/2\); the next two give a summable
\((1+|r|)^{-2}\) interaction between material cells at distance \(r\).
Cauchy--Schwarz after the full cyclic sum yields
\begin{equation}
 |D^2\kform|
 \le Ca\sum_i\gamma_i^2d_i^2.
\end{equation}
The factor \(a\) is the order-minus-one gain.  This is the source-only
specialization of the real-material Bessel Hessian estimate proved in
stage 65; no Galerkin or support-space term is present here.
\end{proof}

\begin{proposition}[Relative Bessel source row]\label{F129:prop:krelative}
For every sufficiently fine positive fan,
\begin{equation}\label{F129:eq:krelative}
 -\{\kform(\PiNG B_\alpha)-\kform(\PiNG B_{a\mathbf1})\}
 \le C\Jfour(\alpha).
\end{equation}
\end{proposition}

\begin{proof}
If \(\Jfour(\alpha)\ge a^5\), positivity and
Lemma~\ref{F129:lem:kbaseline} give
\begin{equation}
 -\{\kform(B_\alpha)-\kform(B_{a\mathbf1})\}
 \le\kform(B_{a\mathbf1})\le C\Jfour(\alpha).
\end{equation}
If \(\Jfour(\alpha)<a^5\), \eqref{F129:eq:Jidentity} implies
\(\max_i|\alpha_i-a|^4\le\Jfour<a^5\).  For small \(a\), the whole affine
segment \(\gamma_i(t)=a+t(\alpha_i-a)\) lies in
\([a/2,3a/2]\).  Cyclic invariance kills the first derivative at
\(a\mathbf1\).  Lemma~\ref{F129:lem:kHessian}, Taylor's formula, and
\begin{equation}
 \Jfour(\alpha)
 =12\int_0^1(1-t)\sum_i\gamma_i(t)^2(\alpha_i-a)^2\,dt
\end{equation}
prove the stronger bound \(Ca\Jfour\).
\end{proof}

\subsection{Exact source envelope and endpoint forcing}

\begin{theorem}[Unnormalized exact disk-source envelope]\label{F129:thm:B}
For every sufficiently fine positive fan,
\begin{equation}\label{F129:eq:B-envelope}
 \qD(\PiNG B_\alpha)-\qD(\PiNG B_{a\mathbf1})
 \le C\Jfour(\alpha).
\end{equation}
\end{theorem}

\begin{proof}
Add Theorem~\ref{F129:thm:qdot}, Lemma~\ref{F129:lem:qzero}, and
Proposition~\ref{F129:prop:krelative} according to the exact multiplier split
\eqref{F129:eq:split}.
\end{proof}

\begin{theorem}[Area-normalized exact source envelope]\label{F129:thm:b}
For every sufficiently fine positive fan,
\begin{equation}\label{F129:eq:b-envelope}
 \boxed{\qD(b_\alpha)-\qD(b_{a\mathbf1})
 \le C\Jfour(\alpha).}
\end{equation}
\end{theorem}

\begin{proof}
Convexity of \(\tan(x/2)\) and \(\sum_i\alpha_i=Na\) give
\(A_\alpha\ge A_{a\mathbf1}\), hence
\(h_\alpha^2\le h_{a\mathbf1}^2\).  From \eqref{F129:eq:b},
\begin{align}
 \qD(b_\alpha)-\qD(b_{a\mathbf1})
 &=h_\alpha^2\{\qD(B_\alpha)-\qD(B_{a\mathbf1})\}\notag\\
 &\quad +(h_\alpha^2-h_{a\mathbf1}^2)\qD(B_{a\mathbf1}).\label{F129:eq:area-sign}
\end{align}
The second term is nonpositive.  The factors \(h_\alpha\) are uniformly
bounded for fine fans, so Theorem~\ref{F129:thm:B} proves the claim.
\end{proof}
\endgroup
